\documentclass[12pt]{article}
\usepackage[utf8]{inputenc}
\usepackage[T1]{fontenc}
\usepackage[margin=1in]{geometry}
\usepackage{setspace}
\usepackage{xcolor}
\usepackage[colorlinks=true, urlcolor=blue!60!black, linkcolor=blue!60!black]{hyperref}
\setlength{\parindent}{0em}
\setlength{\parskip}{0.45em}
\renewcommand{\baselinestretch}{1.0}
\begin{document}
\begin{center}
{\Large \textit{NFIB v. Sebelius}, 567 U.S. 519 (2012)\par}
\vspace{0.3em}
{\normalsize Case Brief --- One-Page Edition\par}
\end{center}

\textbf{Memory Jogger.} The individual mandate exceeds the commerce power but is a valid tax; the Medicaid ultimatum is coercive (\href{https://supreme.justia.com/cases/federal/us/567/519/}{Justia}).

\textbf{Facts.} The Affordable Care Act required most people to carry insurance or pay the IRS a ``shared responsibility payment,'' and required States to expand Medicaid or risk all their Medicaid funds.

\textbf{Procedural History.} The Eleventh Circuit held the mandate unconstitutional but severable.

\textbf{Votes.} Shifting majorities led by Roberts: 5--4 against the commerce theory, 5--4 upholding the tax, 7--2 on coercion.

\textbf{Holding.} A modest, IRS-collected exaction without scienter that leaves a lawful choice is a tax. Threatening existing grants to force a new program is coercion.

\textbf{Reasoning.} Regulating commerce presupposes activity; compelling purchases would create a general police power. But reading the payment as a tax was ``fairly possible.'' The Medicaid threat, over 10\% of state budgets, was ``economic dragooning.'' Ginsburg would have upheld the mandate under the commerce power; the joint dissent would have struck the whole Act.

\textbf{Examples.} \emph{Same-Side:} a modest IRS flood-insurance assessment. \emph{Opposite-Side:} a jail-backed electric-vehicle mandate. \emph{Fence-Sitter:} a steep vaccine assessment.

\textbf{Critique.} Critics call the activity/inactivity line arbitrary and the tax reading textually strained.

\textbf{Key Quotation.} ``The Framers gave Congress the power to regulate commerce, not to compel it, and for over 200 years both our decisions and Congress's actions have reflected this understanding. There is no reason to depart from that understanding now.'' --- Chief Justice Roberts, 567 U.S. at 555.

\end{document}
